\section{Conclusion}
\label{sec:conclusion}

We presented the first single-protocol, head-to-head ablation of six LLM ideation
regimes---one ungrounded control and five grounding mechanisms---holding the model,
task suite, generation budget, and evaluation procedure fixed and scoring each by
objective downstream utility. Our answer to the motivating question is nuanced and
partly negative: under a fixed, objective protocol, \emph{grounding does not reliably
make a strong LLM's ideas better}. Pooled across mechanisms, grounded banks are
statistically indistinguishable from a strong ungrounded control.

\para{Key findings.} The mechanism is what matters. Grounding in empirical feedback
(test-and-revise) is the only regime that significantly improves in-distribution
usefulness ($+4.4$ points, $d{=}1.09$), but it costs $\sim\!80\times$ the tokens and
its gains do not transfer out-of-distribution. Pure elaboration (proposal writing)
significantly \emph{hurts} and collapses idea diversity. Data grounding can
\emph{overfit}, trading out-of-distribution robustness for in-distribution fit. And
novelty and diversity do not track usefulness within a task. Ungrounded prompting of
a capable model is a surprisingly hard baseline to beat.

\para{Future work.} Four directions follow directly. (1)~A cross-model sweep from
weaker to frontier generators, to map where grounding's value lives as the ungrounded
baseline strengthens. (2)~Regularized or \OOD-aware empirical grounding that keeps the
in-distribution gain without overfitting to surface cues. (3)~Idea-level domains
(\eg rediscovery benchmarks) to test whether these conclusions hold for open-ended
ideation rather than classification rules. (4)~Hybrid ``prior $+$ minimal empirical
check'' regimes that aim to capture most of empirical grounding's gain at a fraction
of its cost. We release all banks, per-item predictions, and statistics to support
these next steps.
